\section{Curve25519 and its discrete logarithm}

Bernstein designed Curve25519 for fast Diffie--Hellman and straightforward
constant-time implementations \cite{bernstein2006curve25519}. RFC~7748 gives
the modern interoperable specification \cite{rfc7748}.

\subsection{Parameters}

\begin{center}
\begin{tabularx}{\textwidth}{lX}
\toprule
Field prime & $p=2^{255}-19$ \\
Curve & $v^2=u^3+486662u^2+u$ over $\F_p$ \\
Prime subgroup order &
$\ell=2^{252}+27742317777372353535851937790883648493$ \\
Cofactor & $h=8$, so $\#E(\F_p)=8\ell$ \\
Base point & Montgomery coordinate $u(G)=9$, with $\ord(G)=\ell$ \\
\bottomrule
\end{tabularx}
\end{center}

The name ``25519'' comes from $2^{255}-19$, not from a catalog number or a
claim of 255-bit security. The pseudo-Mersenne shape gives
\[
  2^{255}\equiv19\pmod p,
\]
so high limbs can be folded back into the field efficiently. The public
coefficient $486662$ is not a key: it is part of a parameter choice that gives
useful group and twist orders together with efficient Montgomery formulas.

Since $\ell\approx2^{252}$, a generic Pollard-rho attack costs on the order of
\[
  \sqrt\ell\approx2^{126}
\]
group operations. The curve is therefore conventionally placed near the
128-bit classical security level. Field size and security level are different
quantities.

\subsection{What changes when only the u-coordinate is exposed?}

On a Montgomery curve,
\[
  -(u,v)=(u,-v).
\]
Thus $u(P)=u(-P)$: retaining only $u$ identifies a point with its inverse. For
the odd prime-order subgroup $\langle G\rangle$, this is the only ambiguity.

\begin{proposition}
For nonidentity points in $\langle G\rangle$,
\[
  u([a]G)=u([b]G)
  \quad\Longleftrightarrow\quad
  b\equiv a\ \text{or}\ b\equiv-a\pmod\ell.
\]
\end{proposition}

\begin{proof}
Two nonidentity curve points with the same $u$ have $v$-coordinates that are
negatives, so they are equal or inverses. Hence $[b]G=[a]G$ or
$[b]G=-[a]G$. Since $G$ has order $\ell$, the corresponding scalar
congruences follow. The reverse implication is immediate.
\end{proof}

Consequently, the x-only discrete-log relation for an honest public key is
recovery of $a$ up to sign modulo $\ell$. That is enough to break
Diffie--Hellman, because
\[
  u([-a]B)=u(-[a]B)=u([a]B).
\]
An ECDLP solution recovers a scalar residue, not necessarily the original
32-byte secret input: X25519 deliberately overwrites five input bits during
scalar decoding, and scalar multiplication depends only on the residue in the
subgroup.

The corresponding x-only CDH problem is
\[
  \bigl(u(G),u([a]G),u([b]G)\bigr)
  \longmapsto u([ab]G).
\]
This is the formal connection between the abstract DLP and honest Curve25519
public keys. It does not establish an equivalence between ECDLP and CDH; it
shows that an ECDLP solver is sufficient to solve the Curve25519 CDH instance.

\subsection{The quadratic twist and cofactor}

For a nonsquare $d\in\F_p$, a quadratic twist can be represented as
\[
  E^d:\quad d v^2=u^3+Au^2+u.
\]
The curve and twist orders satisfy
\[
  \#E(\F_p)+\#E^d(\F_p)=2(p+1).
\]
A received $u$ may correspond to the main curve or to its twist. Curve25519's
twist has cofactor $4$ and a large prime factor, so the x-only ladder can be
defined robustly on both sides. Inputs in a small-order component can still
produce the all-zero output. Making the scalar divisible by $8$ annihilates
the curve's cofactor-$8$ component (and the twist's cofactor-$4$ component),
but it does not make every peer input contributory or authenticated.

\begin{takeaway}
The hard group problem lives in the subgroup of order $\ell$, not in all
$2^{255}$ possible byte strings and not in the cofactor-$8$ ambient group.
\end{takeaway}
